% Generated from Review/corrected_text.md (the reviewed source text).

\bakhsh{ج. نیمه‌نهایی‌ها و فینال}

\band{۱. بهترین نتیجه از سه بازی}

بازیکنی که در فهرست نهایی رده‌بندی فردی رتبه بهتری دارد، رنگ بازی اول را انتخاب می‌کند. در هر مسابقه، به هر بازیکن برای هر برد یک امتیاز و برای هر تساوی نیم امتیاز داده می‌شود. رنگ‌ها در بازی دوم جابه‌جا می‌شوند. اگر بازیکنی هر دو بازی اول و دوم را ببرد، برنده اعلام می‌شود و بازی سوم انجام نمی‌شود. حق انتخاب رنگ در بازی سوم (در صورت نیاز) به بازیکنی داده می‌شود که در دو بازی اول تعداد مهره‌های بیشتری داشته است. در صورت تساوی، بازیکنی که در فهرست نهایی رده‌بندی فردی رتبه بهتری دارد، حق انتخاب رنگ را خواهد داشت. در پایان سه بازی، بازیکنی که امتیاز بیشتری در مسابقه کسب کرده باشد، برنده اعلام می‌شود. اگر پس از سه بازی هر دو بازیکن امتیاز برابر داشته باشند، مجموع تعداد مهره‌ها در سه بازی برای تعیین برنده به کار می‌رود. در صورت ادامه تساوی، بازیکنی که در فهرست نهایی رده‌بندی فردی رتبه بهتری دارد، برنده اعلام می‌شود.

\band{۲. مسابقات جهانی اتللو}

\zirband{\LRE{۲\kern0pt.\kern0pt۱} نیمه‌نهایی‌ها}

در آغاز روز سوم رقابت‌ها، چهار بازیکن برتر فهرست نهایی رده‌بندی فردی در نیمه‌نهایی‌ها شرکت خواهند کرد. هر مسابقه نیمه‌نهایی به شکل بهترین نتیجه از سه بازی خواهد بود. نفر اول فهرست نهایی رده‌بندی فردی با نفر چهارم، و نفر دوم با نفر سوم بازی خواهد کرد. برندگان دو مسابقه نیمه‌نهایی به فینال راه می‌یابند.

\zirband{\LRE{۲\kern0pt.\kern0pt۲} پلی‌آف مقام سوم}

دو بازنده نیمه‌نهایی‌ها یک بازی پلی‌آف برای مقام سوم انجام خواهند داد. این بازی طوری زمان‌بندی می‌شود که هم‌زمان با بازی اول فینال مسابقات جهانی اتللو آغاز شود. بازیکنی که در فهرست نهایی رده‌بندی فردی رتبه بهتری دارد، حق انتخاب رنگ یا تساوی-برد را دارد (بند د.۹ را ببینید).

\zirband{\LRE{۲\kern0pt.\kern0pt۳} فینال}

فینال به شکل بهترین نتیجه از سه بازی خواهد بود. به برنده، عنوان قهرمان مسابقات جهانی اتللو اعطا می‌شود.

\band{۳. مسابقات جهانی اتللو زنان}

\zirband{\LRE{۳\kern0pt.\kern0pt۱} نیمه‌نهایی زنان}

اگر یک یا چند زن به نیمه‌نهایی مسابقات جهانی اتللو راه یافته باشند، نیمه‌نهایی زنان برگزار نخواهد شد. در عوض، تمام زنانی که در نیمه‌نهایی مسابقات جهانی اتللو حضور دارند، بر اساس ترتیب نهایی خود پس از فینال مسابقات جهانی اتللو و پلی‌آف مقام سوم رده‌بندی خواهند شد. زنی که رتبه بهتری کسب کند، عنوان قهرمان زنان جهان را دریافت می‌کند. اگر هیچ زنی به نیمه‌نهایی مسابقات جهانی اتللو راه نیافته باشد، نفر اول فهرست نهایی رده‌بندی زنان با نفر چهارم، و نفر دوم با نفر سوم بازی خواهد کرد. هر مسابقه نیمه‌نهایی به شکل تک‌بازی خواهد بود. در هر بازی نیمه‌نهایی، زنی که در فهرست نهایی رده‌بندی زنان رتبه بهتری دارد، حق انتخاب رنگ یا تساوی-برد را خواهد داشت (بند د.۹ را ببینید). برندگان دو نیمه‌نهایی به فینال راه می‌یابند. هر دو بازی طوری زمان‌بندی می‌شوند که هم‌زمان با بازی‌های دوم نیمه‌نهایی مسابقات جهانی اتللو آغاز شوند.

\zirband{\LRE{۳\kern0pt.\kern0pt۲} پلی‌آف مقام سوم زنان}

دو بازنده نیمه‌نهایی‌های زنان یک بازی پلی‌آف برای مقام سوم انجام خواهند داد. این بازی طوری زمان‌بندی می‌شود که هم‌زمان با بازی اول فینال مسابقات جهانی اتللو (و پلی‌آف مقام سوم مسابقات جهانی اتللو) آغاز شود. زنی که در فهرست نهایی رده‌بندی زنان رتبه بهتری دارد، حق انتخاب رنگ یا تساوی-برد را خواهد داشت (بند د.۹ را ببینید).

\zirband{\LRE{۳\kern0pt.\kern0pt۳} فینال زنان}

فینال به شکل یک تک‌بازی بین برندگان دو نیمه‌نهایی زنان برگزار خواهد شد. این بازی طوری زمان‌بندی می‌شود که هم‌زمان با بازی دوم فینال مسابقات جهانی اتللو آغاز شود. زنی که در فهرست نهایی رده‌بندی زنان رتبه بهتری دارد، حق انتخاب رنگ یا تساوی-برد را خواهد داشت (بند د.۹ را ببینید). برنده، عنوان قهرمان زنان جهان را دریافت می‌کند.

\band{۴. مسابقات جهانی اتللو نوجوانان}

\zirband{\LRE{۴\kern0pt.\kern0pt۱} نیمه‌نهایی نوجوانان}

اگر یک یا چند بازیکن نوجوان به نیمه‌نهایی مسابقات جهانی اتللو راه یافته باشند، نیمه‌نهایی نوجوانان برگزار نخواهد شد. در عوض، تمام بازیکنان نوجوانی که در نیمه‌نهایی مسابقات جهانی اتللو حضور دارند، بر اساس ترتیب نهایی خود پس از فینال مسابقات جهانی اتللو و پلی‌آف مقام سوم رده‌بندی خواهند شد. بازیکن نوجوانی که رتبه بهتری کسب کند، عنوان قهرمان نوجوانان جهان را دریافت می‌کند. اگر هیچ بازیکن نوجوانی به نیمه‌نهایی مسابقات جهانی اتللو راه نیافته باشد، نفر اول فهرست نهایی رده‌بندی نوجوانان با نفر چهارم، و نفر دوم با نفر سوم بازی خواهد کرد. هر مسابقه نیمه‌نهایی به شکل تک‌بازی خواهد بود. در هر بازی نیمه‌نهایی، بازیکن نوجوانی که در فهرست نهایی رده‌بندی نوجوانان رتبه بهتری دارد، حق انتخاب رنگ یا تساوی-برد را خواهد داشت (بند د.۹ را ببینید). برندگان دو نیمه‌نهایی به فینال راه می‌یابند. هر دو بازی طوری زمان‌بندی می‌شوند که هم‌زمان با بازی‌های اول نیمه‌نهایی مسابقات جهانی اتللو آغاز شوند.

\zirband{\LRE{۴\kern0pt.\kern0pt۲} پلی‌آف مقام سوم نوجوانان}

دو بازنده نیمه‌نهایی‌های نوجوانان یک بازی پلی‌آف برای مقام سوم انجام خواهند داد. این بازی طوری زمان‌بندی می‌شود که در زمان تعیین‌شده برای بازی‌های سوم نیمه‌نهایی مسابقات جهانی اتللو آغاز شود (حتی اگر بازی سوم در نیمه‌نهایی‌های مسابقات جهانی اتللو لازم نباشد). بازیکن نوجوانی که در فهرست نهایی رده‌بندی نوجوانان رتبه بهتری دارد، حق انتخاب رنگ یا تساوی-برد را خواهد داشت (بند د.۹ را ببینید).

\zirband{\LRE{۴\kern0pt.\kern0pt۳} فینال نوجوانان}

فینال به شکل یک تک‌بازی بین برندگان دو نیمه‌نهایی نوجوانان برگزار خواهد شد. این بازی طوری زمان‌بندی می‌شود که در زمان تعیین‌شده برای بازی‌های سوم نیمه‌نهایی مسابقات جهانی اتللو آغاز شود (حتی اگر بازی سوم در نیمه‌نهایی‌های مسابقات جهانی اتللو لازم نباشد). بازیکن نوجوانی که در فهرست نهایی رده‌بندی نوجوانان رتبه بهتری دارد، حق انتخاب رنگ یا تساوی-برد را خواهد داشت (بند د.۹ را ببینید). برنده، عنوان قهرمان نوجوانان جهان را دریافت می‌کند.
